% Libro 3 · Mundo 5 · Referencia y cohesión (app: pu-referencia)
\mundo{Referencia y cohesión}{pu-referencia}{Un texto no repite todo: dice «ella», «este fenómeno», «lo anterior», o
simplemente calla lo que ya se sabe. Esas piezas mantienen unido el texto y, a la vez, son fuente de preguntas: ¿a quién se
refiere «ella»?, ¿qué es «lo anterior»? Este mundo enseña a seguir los hilos.}{%
\item \textbf{5.1} Pronombres y referentes
\item \textbf{5.2} Sustituciones y elipsis}

\leccion{pu-rf-1}{Pronombres y referentes}{Juego de unir cada pronombre con su referente.}

\tecnica{La técnica: volver atrás y probar}
Un \clave{pronombre} (él, ella, lo, la, les, su, este, aquel, que, cuyo…) apunta a algo mencionado antes, su
\clave{referente}. Para encontrarlo:
\begin{enumerate}
  \item Mira hacia atrás (a veces hacia adelante) y lista los candidatos que concuerdan en género y número.
  \item \emph{Reemplaza} el pronombre por cada candidato y lee la oración.
  \item Elige el que da sentido coherente con el resto del texto.
\end{enumerate}

Los demostrativos tienen además una pista de distancia: \emph{este} suele referirse a lo más cercano y \emph{aquel} a lo más
lejano. Y los neutros (\emph{esto}, \emph{eso}, \emph{lo anterior}) suelen referirse a una idea completa, no a una palabra.

\figura{referente}{El pronombre señala hacia atrás, a su referente.}

\begin{regla}
\begin{itemize}
  \item Concordancia: género y número descartan candidatos.
  \item Reemplazo: lee la oración con cada candidato.
  \item Neutros (\emph{esto}, \emph{ello}, \emph{lo anterior}): buscan una idea entera.
\end{itemize}
\end{regla}

\begin{texto}{El puente de los micos}{Crónica · texto original escrito para este libro}
\pa En la vía que atraviesa el bosque, los atropellos de micos eran frecuentes. Los animales cruzaban por el suelo porque
la carretera había partido en dos la línea de árboles que \emph{ellos} usaban para moverse.

\pa Un grupo de biólogos propuso una solución sencilla: tender cuerdas gruesas entre los árboles de ambos lados, a unos
diez metros de altura. \emph{Esta} idea fue recibida con escepticismo por los ingenieros de la vía, \emph{quienes} dudaban de
que los micos usaran un puente que no conocían.

\pa Los primeros meses les dieron la razón: las cuerdas permanecían vacías. Pero un día, una cámara instalada en un árbol
registró a una hembra cruzando con su cría en la espalda. Después de \emph{aquello}, el uso se volvió diario.

\pa Hoy la carretera tiene siete puentes de cuerda, y los atropellos disminuyeron de manera notable. \emph{Su} costo fue
menor que el de una sola señal de tránsito luminosa.
\end{texto}

\subsection{Preguntas comentadas}
\begin{comentada}{En el \S\,2, \emph{quienes} se refiere a}
\mal{A}{los biólogos.}{Fueron ellos quienes propusieron la idea; no tiene sentido que dudaran de ella.}
\bien{B}{los ingenieros de la vía.}{Es el antecedente inmediato y el que encaja: dudaban del puente.}
\mal{C}{los micos.}{Los micos no «dudaban».}
\mal{D}{los árboles.}{No concuerda en sentido.}
\end{comentada}

\begin{reto}
«María le prestó el libro a Sandra porque ella ya lo había leído.» ¿A quién se refiere, más probablemente, \emph{ella}?
\opcionesd{a María}{a Sandra}{al libro}{a las dos}
\end{reto}

\begin{retoresuelto}
Si «ya lo había leído», tiene sentido que preste el libro quien ya lo leyó: \emph{ella} es María. Con Sandra la causa no
tendría lógica (¿por qué prestarle un libro que ya leyó?). \textbf{Opción A.}
\end{retoresuelto}

\taller
\piso{Como en el examen}{opción múltiple}
\begin{ejercicios}
\item En el \S\,1, \emph{ellos} se refiere a
\opcionesd{los árboles}{los micos}{los conductores}{los biólogos}
\sol{B.}
\item En el \S\,2, \emph{Esta} idea es
\opcionesl{atropellar menos micos.}{tender cuerdas entre los árboles de ambos lados.}{instalar señales de tránsito.}{cerrar la carretera.}
\sol{B.}
\item En el \S\,3, \emph{aquello} se refiere a
\opcionesl{que las cuerdas estuvieran vacías.}{el registro de la hembra cruzando con su cría.}{la instalación de la cámara.}{los primeros meses.}
\sol{B.}
\item En el \S\,4, \emph{Su} costo es el de
\opcionesd{los siete puentes de cuerda}{la carretera}{la señal luminosa}{los atropellos}
\sol{A.}
\end{ejercicios}

\bitacora

\leccion{pu-rf-2}{Sustituciones y elipsis}{Juego de completar lo que el texto calla.}

\tecnica{La técnica: reponer lo que falta}
Para no repetir, los textos usan dos recursos además de los pronombres:
\begin{itemize}
  \item \clave{Sustitución}: se cambia una expresión por otra equivalente. «Gabriel García Márquez» se vuelve luego «el
  autor de \emph{Cien años de soledad}», «el Nobel colombiano», «el escritor».
  \item \clave{Elipsis}: se omite algo que se sobreentiende. «Ana estudia medicina; Pedro, derecho» (se omite
  «estudia»).
\end{itemize}
Las preguntas piden reconocer a qué se refiere la sustitución o qué se omitió. El método es \emph{reponer} lo que falta y
comprobar que la oración tiene sentido.

\begin{regla}
\begin{itemize}
  \item Una expresión nueva («el fenómeno», «la medida») puede estar reemplazando algo dicho antes: búscalo.
  \item Ante una coma que «suena a hueco», repón el verbo de la oración anterior.
\end{itemize}
\end{regla}

\begin{texto}{Dos maneras de contar}{Ensayo breve · texto original escrito para este libro}
\pa Los antiguos mayas contaban de veinte en veinte; los romanos, de diez en diez, aunque con letras. La diferencia no era un capricho: probablemente los primeros usaban los dedos de las manos y de los pies; los segundos, solo los de las manos.

\pa El sistema maya tenía, además, un símbolo para el cero; el romano no lo tenía, y en Europa esa idea tardó siglos en aceptarse. \emph{Esa ausencia}, junto con un sistema que no era posicional, hizo que las cuentas romanas fueran lentas y difíciles de escribir.

\pa Hoy usamos números de origen indio, transmitidos por los árabes. \emph{Este viaje de siglos} explica por qué los llamamos
«arábigos» aunque no hayan nacido en Arabia.
\end{texto}

\subsection{Preguntas comentadas}
\begin{comentada}{En el \S\,2, «Esa ausencia» se refiere a}
\mal{A}{la falta de dedos.}{\emph{Dice otra cosa}.}
\bien{B}{la falta de un símbolo para el cero en el sistema romano (europeo).}{Antes se dice que en Europa el cero tardó
siglos en aceptarse; la consecuencia (cuentas romanas lentas) lo confirma.}
\mal{C}{la ausencia de los mayas en Europa.}{No se habla de eso.}
\mal{D}{la falta de letras en el sistema maya.}{\emph{No lo dice}.}
\end{comentada}

\begin{reto}
«Los estudiantes de la mañana prefieren el examen escrito; los de la tarde, el oral.» ¿Qué verbo se omite después de
«los de la tarde»?
\opcionesd{prefieren}{estudian}{presentan}{escriben}
\end{reto}

\begin{retoresuelto}
Completa: «los [estudiantes] de la tarde [prefieren] el [examen] oral». La coma marca el hueco del verbo de la oración
anterior: \emph{prefieren}. También se sobreentienden «estudiantes» y «examen». \textbf{Opción A.}
\end{retoresuelto}

\taller
\piso{Como en el examen}{opción múltiple}
\begin{ejercicios}
\item En el \S\,1, después de «los romanos» se omite
\opcionesd{contaban}{usaban}{eran}{escribían}
\sol{A. «Los romanos contaban de diez en diez».}
\item En el \S\,1, «los segundos» son
\opcionesd{los mayas}{los romanos}{los dedos}{los números}
\sol{B.}
\item «Este viaje de siglos» \pf{3} se refiere a
\opcionesl{el paso de los números desde la India, a través de los árabes, hasta nosotros.}{un viaje de los mayas a Arabia.}{la historia del cero en Europa.}{el sistema romano.}
\sol{A.}
\end{ejercicios}

\bitacora

\evaluacion
\begin{texto}{La maestra del barco}{Relato · texto original escrito para este libro}
\pa Durante quince años, Teresa dio clases en un barco que recorría los caseríos del río. Los niños la esperaban en el
muelle con los cuadernos bajo el brazo; \emph{ella} bajaba con una caja de libros y una tiza que siempre parecía la última.

\pa El barco paraba dos días en cada pueblo. En ese tiempo, la maestra enseñaba a leer a los pequeños; a los grandes,
cuentas y geografía. \emph{Lo poco que alcanzaba a enseñar} lo dejaba anotado en una libreta para que los padres continuaran.

\pa Un año el río se secó tanto que el barco no pudo pasar. Teresa no se resignó: consiguió una mula y siguió el
recorrido por la orilla. \emph{Aquella} temporada tardó el doble, pero no faltó a ningún pueblo.

\pa Cuando se jubiló, los antiguos alumnos le regalaron un mapa del río con los nombres de todos los niños que \emph{la}
habían esperado en el muelle.
\end{texto}
\begin{ejercicios}
\item En el \S\,1, \emph{ella} se refiere a
\opcionesd{la tiza}{Teresa}{la caja}{la libreta}
\sol{B. (5.1)}
\item En el \S\,2, después de «a los grandes» se omite
\opcionesd{enseñaba}{leía}{esperaba}{dejaba}
\sol{A. (5.2)}
\item «Lo poco que alcanzaba a enseñar» \pf{2} se refiere a
\opcionesl{lo que enseñaba en los dos días de cada parada.}{las cuentas de los adultos.}{los libros de la caja.}{las clases de los padres.}
\sol{A. (5.2)}
\item En el \S\,3, \emph{Aquella} temporada es
\opcionesl{la del año en que el río se secó.}{la primera de Teresa.}{la de su jubilación.}{la de las lluvias.}
\sol{A. (5.1)}
\item En el \S\,4, \emph{la} se refiere a
\opcionesd{la mula}{Teresa}{la libreta}{la orilla}
\sol{B. (5.1)}
\end{ejercicios}
